% MixDom-specific extension of the SVI-BW convergence appendix.
% Loaded only by mixdom.tex; the TKF-only paper omits it.
\section{MixDom-Specific SVI-BW Convergence Considerations}
\label{sec:svb-convergence-mixdom}

This appendix specialises the model-agnostic convergence analysis
of Appendix~\ref{sec:svb-convergence} and the BDI expected-statistic
formulae of Appendix~\ref{sec:expected-stats} to the hierarchical
MixDom model.  Each subsection corresponds to a question that arises
specifically because MixDom carries multiple interacting parameter
groups (top-level vs.\ per-domain BDI, intra-fragment fragment-type
chains, and per-class substitution models).

\subsection{Parameter groups and Fisher information}
\label{sec:svb-mixdom-param-groups}

The MixDom model has several parameter groups, each with different
Fisher information characteristics:

\begin{center}
\begin{tabular}{lccc}
\toprule
Parameter group & \# params & Info per pair & Bottleneck \\
\midrule
Top-level $(\insrate_0, \delrate_0)$ & 2 & $O(1)$ & few domains per sequence \\
Per-domain $(\insrate_d, \delrate_d)$ & $2N_{\text{dom}}$ & $O(w_d)$ & domain frequency $w_d$ \\
Intra-fragment fragment-type transitions $\ext^{(d)}_{fg}$ & $N_{\text{dom}} F^2$ & $O(w_d \bar L_d)$ & fragment count \\
Domain weights $w_d$ & $N_{\text{dom}}-1$ & $O(1)$ & multinomial \\
Substitution $Q_c$ & $O(|\mathcal{A}|^2 C)$ & $O(\bar L \cdot w_c)$ & alignment length \\
\bottomrule
\end{tabular}
\end{center}

The bottleneck for SVB convergence is the \textbf{indel parameters of rare
domains}. A domain with frequency $w_d = 0.05$ contributes useful
BDI statistics to only $\sim 5\%$ of pairs, so its effective sample
size is $\sim 0.05 B$ per minibatch.

\subsection{Substitution vs.\ indel information}
\label{sec:svb-mixdom-subst-vs-indel}

Substitution parameters benefit from $O(\bar L)$ information per pair
(one observation per aligned column), while indel parameters get $O(1)$
information per pair.
For a protein with $\bar L \approx 200$ residues, substitution parameters
converge $\sim 200\times$ faster than indel parameters.
This motivates \textbf{decoupled update frequencies}: update substitution
parameters every minibatch, but average indel parameter estimates over
multiple minibatches.

\subsection{MixDom expected statistics}
\label{sec:svb-mixdom-expected-stats}

We now aggregate the per-process expectations of
Appendix~\ref{sec:expected-stats} to the full MixDom model
with $N_{\text{dom}}$ domain types.

\subsubsection{Top-level (domain birth-death)}

The top-level BDI process has rates $(\insrate_0, \delrate_0)$ with
$\kappa_0 = \insrate_0/\delrate_0$ and creates/destroys domains.
At stationarity:
\begin{alignat}{2}
L_0 &= \kappa_0/(1-\kappa_0)
  &\quad& \text{(expected \# domains per sequence)}
\label{eq:L0}
\\
\expect[B_0] &= \insrate_0\,\evoltime/(1-\kappa_0)
  &\quad& \text{(domain births per pair)}
\label{eq:EB0}
\\
\expect[D_0] &= \expect[B_0]
  &\quad& \text{(domain deaths per pair)}
\label{eq:ED0}
\\
\expect[S_0] &= L_0\,\evoltime
  &\quad& \text{(time-integrated domain count)}
\label{eq:ES0}
\\
M_0 &= 1, \quad T_0 = \evoltime
  &\quad& \text{(one endpoint, one process)}
\label{eq:M0T0}
\end{alignat}

\subsubsection{\texorpdfstring{Per-domain type $d$ (fragment birth-death within a domain)}{Per-domain type d (fragment birth-death within a domain)}}

Each surviving domain link of type~$d$ (with probability $w_d$) contains
a TKF92 fragment process with rates $(\insrate_d, \delrate_d)$,
$\kappa_d = \insrate_d/\delrate_d$.

The number of domain links of type~$d$ is approximately
$L_d = w_d\,L_0$ at stationarity.
Each such link contributes one independent BDI fragment process,
so the \emph{aggregated} fragment-level statistics for domain type~$d$
are:
\begin{alignat}{2}
L_d &= w_d\,L_0\;\cdot\;\frac{\kappa_d}{1-\kappa_d}
  &\quad& \text{(total fragment links, type $d$)}
\label{eq:Ld}
\\
\expect[B_d] &= w_d\,L_0\;\cdot\;
  \frac{\insrate_d\,\evoltime}{1-\kappa_d}
  &\quad& \text{(fragment births, type $d$, per pair)}
\label{eq:EBd}
\\
\expect[D_d] &= \expect[B_d]
  &\quad& \text{(fragment deaths, type $d$, per pair)}
\label{eq:EDd}
\\
\expect[S_d] &= w_d\,L_0\;\cdot\;
  \frac{\kappa_d\,\evoltime}{1-\kappa_d}
  &\quad& \text{(time-integrated fragment count, type $d$)}
\label{eq:ESd}
\\
M_d &= w_d\,L_0
  &\quad& \text{(\# independent processes)}
\label{eq:Md}
\\
T_d &= M_d\,\evoltime = w_d\,L_0\,\evoltime
  &\quad& \text{(total observation time)}
\label{eq:Td}
\end{alignat}
Note that $M_d$ and $T_d$ are themselves random (they depend on how
many domains of type~$d$ survive), but at stationarity their
expectations are as above.

\subsubsection{Intra-fragment fragment-type transitions and sequence length}

Within each fragment of domain $d$ the fragment-type chain
is governed by an $F \times F$ Markov transition matrix $\ext^{(d)}_{fg}$
(intra-fragment; different fragments are independent realisations).
The termination probability from fragment-type $f$ is
$\notext^{(d)}_f = 1 - \sum_g \ext^{(d)}_{fg}$.
The expected sojourn length starting from fragment state $f$ is determined
by the fundamental matrix $(I - \ext^{(d)})^{-1}$: specifically, the expected
number of sites emitted starting from a fragment initiated in state $f$ is
$\sum_g [(I - \ext^{(d)})^{-1}]_{fg}$.
Averaging over the initial fragment distribution $\fragdist_{d,f}$, the
expected number of residues per fragment is
$\bar{K}_d = \sum_f \fragdist_{d,f} \sum_g [(I - \ext^{(d)})^{-1}]_{fg}$.

The expected number of residues per domain of type~$d$ is:
\begin{equation}\label{eq:chars-per-domain}
\bar{C}_d = \frac{\kappa_d}{1-\kappa_d}\, \bar{K}_d,
\end{equation}
and the total expected sequence length is:
\begin{equation}\label{eq:total-length}
\bar{L}_{\text{seq}} = L_0 \sum_d w_d\,\bar{C}_d
  = \frac{\kappa_0}{1-\kappa_0}
    \sum_d w_d\,\frac{\kappa_d}{1-\kappa_d}\, \bar{K}_d.
\end{equation}

\begin{remark}[Convergence of Markov chain transition counts]
When fragments are IID (the $\nfrag = 1$ scalar-extension special
case), the sufficient statistics for the fragment extension parameter
$\ext_f$ are Bernoulli counts (extend vs.\ terminate), and the
convergence analysis reduces to a Beta-distributed posterior.
With the general fragment-type Markov, the sufficient statistics are
rows of a Markov chain transition count matrix: for each domain $d$
and source fragment state $f$, we observe
counts $\hat{n}^{(d)}_{fg}$ of transitions to each target state $g$
plus termination counts $\hat{n}^{(d)}_{f,\text{end}}$.
The M-step row-normalizes these counts (a Dirichlet posterior), and the
per-pair Fisher information for each row scales with the expected number
of visits to state $f$ per pair. For fragment states that are rarely
visited (low stationary probability under $\ext^{(d)}$), the effective
sample size is correspondingly small, mirroring the rare-domain bottleneck
at the top level.
\end{remark}


\subsubsection{Numerical example: d3 checkpoint}

We evaluate these formulas using the MixDom d3 checkpoint parameters:
\begin{alignat*}{3}
&\insrate_0 = 0.01328, \quad &&\delrate_0 = 0.01412,
  \quad &&\kappa_0 = 0.9405 \\
&w = (0.662,\; 0.075,\; 0.264) \\
&\insrate_d = (0.00302,\; 0.01033,\; 0.12370) \\
&\delrate_d = (0.00372,\; 0.04686,\; 0.17775) \\
&\kappa_d = (0.812,\; 0.220,\; 0.696) \\
&\evoltime = 0.5
\end{alignat*}

\begin{center}
\renewcommand{\arraystretch}{1.3}
\begin{tabular}{lcccc}
\toprule
& Top-level & Dom 1 & Dom 2 & Dom 3 \\
& ($d{=}0$) & ($d{=}1$) & ($d{=}2$) & ($d{=}3$) \\
\midrule
$\kappa$ & 0.941 & 0.812 & 0.220 & 0.696 \\
$L = \kappa/(1{-}\kappa)$ & 15.81 & 4.31 & 0.283 & 2.29 \\
$w_d \cdot L_0$ & --- & 10.47 & 1.19 & 4.17 \\
$\expect[B]$ per pair & 0.112 & 0.084 & 0.008 & 0.849 \\
$\expect[D]$ per pair & 0.112 & 0.084 & 0.008 & 0.849 \\
$\expect[S]$ & 7.91 & 22.58 & 0.168 & 4.78 \\
$M$ & 1 & 10.47 & 1.19 & 4.17 \\
$T$ & 0.5 & 5.23 & 0.59 & 2.09 \\
$v_\theta$ (approx.) & $\sim 18$ & $\sim 24$ & $\sim 255$ & $\sim 2.4$ \\
\bottomrule
\end{tabular}
\end{center}

\paragraph{Computation of the table entries.}

\emph{Top level:}
$L_0 = 0.9405 / 0.0595 = 15.81$.
$\expect[B_0] = 0.01328 \times 0.5 / 0.0595 = 0.112$.
$\expect[S_0] = 15.81 \times 0.5 = 7.91$.

\emph{Domain type 1} ($w_1 = 0.662$):
$M_1 = 0.662 \times 15.81 = 10.47$ independent fragment processes.
$L_1 = 10.47 \times 0.812/0.188 = 10.47 \times 4.31 = 45.1$ total fragment links.
$\expect[B_1] = 10.47 \times 0.00302 \times 0.5/0.188 = 0.084$.
(The table shows $L$ per process and $\expect[B]$ aggregated over $M_d$ processes.)

\emph{Domain type 3} ($w_3 = 0.264$):
$M_3 = 0.264 \times 15.81 = 4.17$.
$\expect[B_3] = 4.17 \times 0.12370 \times 0.5/0.304 = 0.849$.
This domain has high $\insrate_3$, so fragment births are frequent
and its indel parameters are easy to estimate ($v \approx 2.4$).

$v_{\insrate_0}$: Using~\eqref{eq:v-lambda-approx} with $\rho \approx 0.5$:
$v_{\insrate_0} \approx 2/\expect[B_0] = 2/0.112 = 17.9$.

$v_{\insrate_2}$: $\expect[B_2] = 1.19 \times 0.01033 \times 0.5 / 0.780 = 0.0079$.
$v_{\insrate_2} \approx 2/0.008 = 255$.


\subsection{Convergence rate estimates}
\label{sec:svb-mixdom-rate-estimates}

From~\eqref{eq:minibatch} and the per-pair relative variance
$v_\theta$, the number of pairs $N = BK$ needed for target relative
error $\varepsilon$ is:
\begin{equation}\label{eq:N-needed}
N \;\geq\; \frac{v_\theta}{\varepsilon^2}.
\end{equation}

\begin{center}
\renewcommand{\arraystretch}{1.3}
\begin{tabular}{lccccc}
\toprule
Parameter & $v_\theta$ & $N$ for $\varepsilon{=}10\%$
  & $N$ for $\varepsilon{=}5\%$ & $N$ for $\varepsilon{=}1\%$ \\
\midrule
$\insrate_0$ (top-level ins) & 18 & 1\,800 & 7\,200 & 180\,000 \\
$\delrate_0$ (top-level del) & 18 & 1\,800 & 7\,200 & 180\,000 \\
$\insrate_1$ (dom 1 ins) & 24 & 2\,400 & 9\,600 & 240\,000 \\
$\insrate_2$ (dom 2 ins) & 255 & 25\,500 & 102\,000 & 2\,550\,000 \\
$\insrate_3$ (dom 3 ins) & 2.4 & 240 & 960 & 24\,000 \\
$w_d$ (domain weights) & $\sim 1/L_0 \approx 0.06$ & 6 & 25 & 630 \\
$\ext^{(d)}_{fg}$ (fragment trans.) & $\sim 1/\bar{C}_d$ & \multicolumn{3}{c}{depends on domain} \\
Substitution ($Q$) & $\sim 1/\bar{L}_{\text{seq}}$ & \multicolumn{3}{c}{$\ll 100$} \\
\bottomrule
\end{tabular}
\end{center}

\subsection{Discussion: why top-level indel rates are hardest}

The table reveals a clear hierarchy of estimation difficulty:

\paragraph{1. Substitution parameters are easiest.}
Each aligned residue pair contributes one independent observation
to the substitution sufficient statistics.
With $\bar{L}_{\text{seq}} \approx 200$--$400$ residues per sequence,
the per-pair Fisher information for substitution parameters is
$O(\bar{L}_{\text{seq}})$, so the per-pair relative variance is
$O(1/\bar{L}_{\text{seq}})$.
A few dozen pairs suffice for $\varepsilon = 5\%$ accuracy.

\paragraph{2. Domain weights converge fast.}
Each domain in the ancestor contributes one multinomial observation
of the domain type.
With $L_0 \approx 16$ domains per sequence, the per-pair
Fisher information is $O(L_0)$, giving $v_{w_d} \sim 1/(w_d L_0)$.
For the most common domain type ($w_1 = 0.662$), $v_{w_1} \approx 0.095$,
and even for the rarest ($w_2 = 0.075$), $v_{w_2} \approx 0.84$.
Domain weights are precisely estimated with $N \sim 100$ pairs.

\paragraph{3. Fragment-type transition parameters are moderate.}
Each intra-fragment fragment-type transition contributes one observation
of the $F \times F$ Markov chain row $\ext^{(d)}_{f,:}$.
With $\sim L_0 w_d \kappa_d/(1-\kappa_d)$ fragments of type~$d$ per pair,
the per-pair information scales with the total fragment count.
For domain 1 (which dominates), there are $\sim 45$ fragment links,
so the per-row variance scales as $\sim F/(45 \cdot \pi_f)$ where
$\pi_f$ is the stationary probability of fragment state $f$.
Domain 2 fragments are rare ($\sim 0.3$ per pair) and harder to estimate.

\paragraph{4. Indel parameters for active domains are easy.}
Domain type~3 has high insertion rate ($\insrate_3 = 0.124$)
and $\expect[B_3] \approx 0.85$ births per pair,
giving $v_{\insrate_3} \approx 2.4$.
A few hundred pairs suffice for $5\%$ accuracy.

\paragraph{5. Indel parameters for common but slow domains are moderate.}
Domain type~1 has $\expect[B_1] \approx 0.084$ births per pair,
giving $v_{\insrate_1} \approx 24$.
This requires $N \approx 10{,}000$ pairs for $5\%$ accuracy---feasible
but not trivial.

\paragraph{6. Indel parameters for rare domains are the bottleneck.}
Domain type~2 has weight $w_2 = 0.075$, low $\kappa_2 = 0.22$
(short domains), and low $\insrate_2 = 0.01$.
The expected fragment births per pair are only $\expect[B_2] \approx 0.008$.
This means on average, only 1 in $\sim 125$ pairs shows even a single
fragment birth event in a domain of type~2.
The per-pair relative variance $v_{\insrate_2} \approx 255$ requires
$N > 100{,}000$ pairs for $5\%$ accuracy.

\paragraph{7. Top-level indel rates are intrinsically hard.}
The top-level BDI controls domain creation and destruction.
With $\insrate_0 = 0.013$, $\evoltime = 0.5$, and $L_0 \approx 16$
domains, the expected number of domain births per pair is only
$\expect[B_0] \approx 0.11$.
The per-pair relative variance $v_{\insrate_0} \approx 18$ requires
$N \approx 7{,}200$ pairs for $5\%$ accuracy.

The fundamental reason is that domain-level indel events are rare
compared to residue-level observations.
A sequence of 300 residues organized into 16 domains gives
$\sim 300$ substitution observations but only $\sim 0.1$ domain birth
observations per pair at $\evoltime = 0.5$.
The information ratio is roughly $300 / 0.1 = 3{,}000 \times$ in favor
of substitution parameters.

\paragraph{Implications for training.}
These estimates motivate the decoupled update strategy recommended
in Section~5: substitution parameters can be accurately estimated
from small minibatches ($B \sim 10$), while indel parameters
(especially for rare domain types) require accumulation over
$N \sim 10^4$--$10^5$ pairs.
The Maraschino pipeline (Section~5.5) achieves this by processing
all training pairs in a single count tensor, at the cost of the
composite-likelihood efficiency gap.
The hybrid Maraschino~$\to$~SVB pipeline is optimal: Maraschino
provides accurate initial estimates using all data, and SVB refines
the indel parameters using the full model.
